% Pitch marks above syllables.  A rising pitch is marked with one or more rising
% diagonal strokes and a falling pitch with falling ones.  The strokes are
% \diagup ("/") and \diagdown ("\"), ordinary glyphs set in \markfont, so --
% unlike the \pdfliteral strokes used before -- nothing here depends on the
% output driver, and the file is fully portable.  \risetwo/\risethree and
% \falltwo/\fallthree stack two and three strokes vertically, one above
% another.  \hold{...} (see hold.tex) underlines a held syllable; since it
% sits below the baseline it composes with the marks above, e.g.
% \rise{\hold{heer}}.
%
% \markfont sizes the marks (its \diagup/\diagdown glyphs are the strokes);
% scale it to make the marks bigger or smaller.  \diagup lives at code 8 and
% \diagdown at code 72 of line10 -- the line-segment font LaTeX's own
% \line(dx,dy) picture-mode command draws with, present in every TeX
% installation.  Its characters are straight strokes at fixed slopes rather
% than a symbol-font glyph; code 8 is \line(2,1) (slope 1/2, ~26.6 degrees
% from the horizontal -- shallower than the ~45-degree diagonal used before)
% and code 72 is \line(2,-1), its mirror image falling instead of rising.
% line10 has a 10pt design size.
\font\markfont=line10 scaled 840    % marks: line10 line segments; \diagup = rise, \diagdown = fall
\chardef\diagup=8                   % rising diagonal stroke  ("/"), slope 1/2
\chardef\diagdown=72                % falling diagonal stroke ("\"), slope 1/2
% line10 draws each stroke with a round pen (a fixed perpendicular width,
% proportional to \markfont's scale -- about 0.34pt at the current 840) too
% thin on its own, so it's emboldened the way LaTeX's own \boldsymbol does
% for AMS symbols -- "poor man's bold": overprint the glyph \boldcopies
% times, each copy shifted \markbold further to the right, so the copies
% merge sideways into one heavier stroke.  A horizontal shift thickens a
% line perpendicular to itself by shift*sin(angle from the horizontal), so
% the same \markbold does less work at this shallower ~26.6 degrees than it
% would nearer 45; \markbold is sized up to compensate.
% \markbold is the step between copies: it must stay smaller than the base stroke
% width or the copies show as separate lines instead of merging.  The perpendicular
% thickness added is (\boldcopies-1)*\markbold*sin(26.6deg), so together with the
% thin base glyph the strokes here are tuned to match the \hold underline (1.44pt,
% measured directly off the rendered strokes -- not just calculated, since both
% \markfont's angle and its scale factor into the base glyph's own thickness).
\newdimen\markbold \markbold=0.23pt      % step between overprinted copies
\newcount\boldcopies \boldcopies=12      % number of copies (weight of the stroke)
\newdimen\boldshift                      % scratch: recentring offset, set per mark
% The strokes are drawn in colour.  Colour is inherently a driver feature, so --
% unlike the stroke shapes -- it needs a pdfTeX \pdfliteral: \markcolor sets the
% (non-stroking / fill) colour for the glyphs and \markcolorreset restores black
% for the surrounding text.  These are emitted inside the mark's own \hbox and
% balanced, so the colour never leaks onto the text or the \hold underlines.
% "1 0 0 rg" is red; change the RGB triple to recolour the marks.
\def\markcolor{\pdfliteral{1 0 0 rg}}    % strokes in red
\def\markcolorreset{\pdfliteral{0 g}}    % text back to black
% \riseheight is the fixed distance from the text baseline to a mark's baseline,
% so every mark sits the same distance above the text and they line up
% horizontally (rather than following each syllable's own height).  Its value is
% measured from the body font in layout.tex; reset it (e.g. per paragraph) to
% move all following marks together.
\newdimen\riseheight
% \risemarkalign{leftglue}{rightglue}{strokes}{syllable} is the shared core: it
% measures the syllable's width, then \rlap overlays the stroke characters --
% horizontally placed over the syllable per leftglue/rightglue and raised so
% their baseline sits \riseheight above the text baseline -- without consuming
% horizontal space, before printing the syllable itself.  \risemark (centred)
% and \risemarkleft (flush left) below are the two alignments built on it.
\def\risemarkalign#1#2#3#4{%
  \leavevmode                 % ensure horizontal mode, so a mark at the start
                              % of a line stays inline with the text
  \setbox0=\hbox{#4}%
  \setbox2=\hbox to\wd0{#1\markfont#3#2}%   the strokes, aligned over the syllable
  % Recentre the smear: it spreads (\boldcopies-1) steps to the right, so shift
  % the whole cluster left by half that to keep the marks centred on themselves.
  \boldshift=\markbold \multiply\boldshift by\boldcopies \advance\boldshift by-\markbold
  \divide\boldshift by2
  \rlap{\raise\riseheight\hbox{%     poor-man's bold: overprint \boldcopies copies,
    \markcolor                       % strokes in colour (reset to black below)
    \kern-\boldshift                 % each \markbold further right (recentred);
    \count255=0                      % the enclosing \rlap keeps the marks from
    \loop \ifnum\count255<\boldcopies %  consuming any horizontal space
      \rlap{\copy2}\kern\markbold
      \advance\count255 by1
    \repeat
    \markcolorreset                  % back to black for the text
  }}%
  % No strut needed: room for the marks is provided globally by the fixed
  % \baselineskip set with the font, so all lines keep the same height.
  \box0
}
\def\risemark#1#2{\risemarkalign{\hss}{\hss}{#1}{#2}}      % strokes centred over the syllable
\def\risemarkleft#1#2{\risemarkalign{}{\hss}{#1}{#2}}      % strokes flush with the syllable's left edge
\def\rise{\risemark{\diagup}}
% \risemarkstackalign{leftglue}{rightglue}{stroke}{count}{syllable}: like
% \risemarkalign, but stacks \stroke \count times vertically -- one directly
% above another -- rather than side by side, for \risetwo/\risethree and
% \falltwo/\fallthree.  Rows are separated by \markstackgap; the vbox's own
% reference point sits at the bottom row's baseline, so \raise\riseheight
% below puts that bottom stroke exactly where a single \rise/\fall would
% sit, with the rest stacking upward from it.  Otherwise the same as
% \risemarkalign: poor-man's-bold smear, colour, then the syllable itself.
% \diagup/\diagdown carry several points of empty padding in their own
% bounding box (msbm10's design, not this file's doing), so a small positive
% \markstackgap is invisible against it -- closing the visible gap needs a
% negative value, overlapping the boxes.
\newdimen\markstackgap \markstackgap=-0.75pt   % vertical gap between stacked strokes
\def\risemarkstackalign#1#2#3#4#5{%
  \leavevmode
  \setbox0=\hbox{#5}%
  \setbox2=\vbox{%
    \offinterlineskip          % rows are spaced by \markstackgap alone, not
                                % the document's normal interline glue
    \count255=0
    \loop \ifnum\count255<#4
      \ifnum\count255>0 \kern\markstackgap\fi
      \hbox to\wd0{#1\markfont#3#2}%
      \advance\count255 by1
    \repeat
  }%
  \boldshift=\markbold \multiply\boldshift by\boldcopies \advance\boldshift by-\markbold
  \divide\boldshift by2
  \rlap{\raise\riseheight\hbox{%
    \markcolor
    \kern-\boldshift
    \count255=0
    \loop \ifnum\count255<\boldcopies
      \rlap{\copy2}\kern\markbold
      \advance\count255 by1
    \repeat
    \markcolorreset
  }}%
  \box0
}
\def\risemarkstack#1#2{\risemarkstackalign{\hss}{\hss}{#1}{#2}}
\def\risetwo{\risemarkstack{\diagup}{2}}
\def\risethree{\risemarkstack{\diagup}{3}}
% \fall, \falltwo, \fallthree: the same, with falling diagonals for the falling pitch.
\def\fall{\risemark{\diagdown}}
\def\falltwo{\risemarkstack{\diagdown}{2}}
\def\fallthree{\risemarkstack{\diagdown}{3}}
% "Spaced" marks: several strokes over one syllable, each corresponding to a
% separate note, so they need clear air between them -- more than the plain
% multi-stroke marks above, which just butt the glyphs together.  \markgap is
% that extra kern, inserted between each pair of strokes (not at the ends).
\newdimen\markgap \markgap=3pt
\def\risetwice{\risemark{\diagup\kern\markgap\diagup}}
\def\fallthreespaced{\risemark{\diagdown\kern\markgap\diagdown\kern\markgap\diagdown}}
\def\falltwice{\risemark{\diagdown\kern\markgap\diagdown}}
\def\fallthreerisespacedleft{\risemarkleft{\diagdown\kern\markgap\diagdown\kern\markgap\diagdown\kern\markgap\diagup}}
\def\risethreefallthreespacedleft{\risemarkleft{\diagup\kern\markgap\diagup\kern\markgap\diagup\kern\markgap\diagdown\kern\markgap\diagdown\kern\markgap\diagdown}}
% \level{...} marks a level (unchanging) pitch with a horizontal rule over the
% syllable, positioned so its top edge lines up with the bottom tip of the
% \rise/\fall diagonals: \diagup and \diagdown both carry the same fixed depth
% below their own baseline in \markfont (\markdepth, measured once below), so
% that baseline -- \riseheight above the text baseline, same as the diagonals --
% less \markdepth is the height of their lower tip, and where the rule's top
% goes. The rule is \holdthick thick and \holdgap narrower than the syllable,
% centered, matching \hold's underline (hold.tex) so the two kinds of marks
% read as one family; \level therefore needs hold.tex input alongside it.
\newdimen\markdepth
\setbox1=\hbox{\markfont\diagdown}%        % \diagup has the same depth
\markdepth=\dp1
\def\level#1{%
  \leavevmode
  \setbox0=\hbox{#1}%
  \dimen0=\wd0 \advance\dimen0 by-\holdgap   % shorten the rule by a fixed amount
  \dimen1=\riseheight \advance\dimen1 by-\markdepth  % height of the rule's top
  \rlap{\kern.5\holdgap                       % ... and centre it over the syllable,
    \raise\dimen1\hbox{%
      \markcolor
      \vrule height\holdthick depth0pt width\dimen0
      \markcolorreset
    }}%
  \box0                                        % leaving half the gap on each side
}
% \levelstroke is a level (flat) stroke sized to drop into a \risemark(align)
% strokes list alongside \diagup/\diagdown, e.g. in \levelrisespaced below.
% msbm10 has no horizontal-bar glyph, so it's a \vrule of \strokewidth (the
% same width as \diagup/\diagdown, measured above into \wd1) and \holdthick
% thick, positioned exactly like \level's rule -- \markdepth below the
% strokes' shared baseline, so its bottom edge lines up with the diagonals'
% bottom tips.  Being a box rather than a font character, it does not pick up
% the poor-man's-bold smear \risemarkalign applies to the whole strokes row
% (that trick thickens diagonal/vertical edges via a sideways shift, which
% would only lengthen a horizontal bar, not thicken it); its \holdthick
% height is already the finished weight.
\newdimen\strokewidth \strokewidth=\wd1
\def\levelstroke{\lower\markdepth\hbox{\vrule height\holdthick depth0pt width\strokewidth}}
% \levelrisespaced{...}: one level stroke followed by one rising stroke, each
% a separate note.  The level stroke's own ink reaches further right than a
% diagonal's ("test van" comparison), so plain \markgap reads as a smaller
% gap than the same kern does between two diagonals; \levelgap adds a bit
% more air to compensate.
\newdimen\levelgap \levelgap=\markgap \advance\levelgap by1.5pt
\def\levelrisespaced{\risemark{\levelstroke\kern\levelgap\diagup}}
\def\level{\risemark{\levelstroke}}